\begin{proposition}[Leray-edge vanishing of normalized cusp restrictions]
\hypertarget{tag-03IM}{}\label{wbII:prop:leray-edge-restriction}
\StatementRecord{wbII:prop:leray-edge-restriction}{proved}{03IM}{\SeriesRef{wbII:eq:W:5.17d}, \SeriesRef{wbII:eq:W:5.17e}, \SeriesRef{wbII:prop:pic0-pullback}; \cite[Tags~012O, 012U and~012W]{StacksProject}; \cite[Tag~09V6]{StacksProject}; \cite[Tags~011N, 011O and~012U]{StacksProject}}

For the constructed map $f:X\to B$ and every
$L\in\operatorname{Pic}(X)$, the degree-two Leray edge section satisfies
\[
 e(c_1(L))=0\in H^0(B,R^2f_*\mathbb Z),
 \qquad
 \eta^*i^*L\simeq\mathcal O_{\mathrm{dP}_6}.
\tag{W.20}\hypertarget{tag-03GM}{}\label{wbII:eq:W:5.17c}
\]
\end{proposition}



\begin{proof}

\begin{enumerate}[label=\arabic*.,leftmargin=*]
\item For every $b\in B$,
proper base change and functoriality of the filtered-complex spectral
sequence give a commutative square
\[
\begin{CD}
 H^2(X;\mathbb Z) @>{e}>> H^0(B,R^2f_*\mathbb Z)\\
 @V{i_b^*}VV @VV{\operatorname{ev}_b}V\\
 H^2(X_b;\mathbb Z) @= H^2(X_b;\mathbb Z).
\end{CD}
\tag{W.21}\hypertarget{tag-03GN}{}\label{wbII:eq:W:5.17d}
\]
Indeed, pull the Leray filtered complex back along $\{b\}\to B$; the
spectral sequence over a point has only column zero, so its edge is the
identity, while the right vertical map is the proper-base-change
identification.  This uses the functoriality and convergence statements in
\cite[Tags~012O, 012U and~012W]{StacksProject} and topological proper base
change \cite[Tag~09V6]{StacksProject}.  The connecting homomorphism in the
exponential sequence is natural under $i_b^*$, so \eqref{wbII:eq:W:5.17d} sends
$c_1(L)$ to $c_1(L|_{X_b})$.

\item Put $s_L=e(c_1(L))$.  The image of an edge morphism consists of permanent
cycles, hence
\[
 s_L\in E_\infty^{0,2}\subseteq\ker
 \bigl(d_2^{0,2}:E_2^{0,2}\to E_2^{2,1}\bigr);
\tag{W.22}\hypertarget{tag-03GO}{}\label{wbII:eq:W:5.17e}
\]
see \cite[Tags~011N, 011O and~012U]{StacksProject}.  For the standard family, \cite[Propositions~7.26--7.27, pp.~61--63]{AC2026}
give $E_2^{0,2}=\mathbb Z(2q)$ and
$E_2^{2,1}=\mathbb Z\,\omega_{\rm AC}[\delta]$, with
$d_2^{0,2}(2q)=\varepsilon p\,\omega_{\rm AC}[\delta]$,
where $\varepsilon\in\{1,-1\}$ is the common differential sign and
$p=12\ell_0-4\ell_1-3\ell_2=-1$.
Define the generator $\omega:=\varepsilon p\,\omega_{\rm AC}$; equivalently,
$\omega[\delta]$ is the actual element $d_2^{0,2}(2q)$.  It is primitive
because $\varepsilon p\in\{1,-1\}$.  With this generator,
\[
 E_2^{0,2}=\mathbb Z(2q),\qquad
 E_2^{2,1}=\mathbb Z\,\omega[\delta],\qquad
 d_2^{0,2}(2q)=\omega[\delta].
\]
Thus $d_2^{0,2}$ is an isomorphism and \eqref{wbII:eq:W:5.17e} forces $s_L=0$.
Taking $b=p_0$ in \eqref{wbII:eq:W:5.17d} gives $c_1(i^*L)=0$, and hence
$c_1(\eta^*i^*L)=0$.  Finally, the injectivity of
$c_1:\operatorname{Pic}(\mathrm{dP}_6)\to H^2(\mathrm{dP}_6;\mathbb Z)$
proved in \cref{wbII:prop:pic0-pullback} makes $\eta^*i^*L$ trivial.

\end{enumerate}
\end{proof}
